% Options for packages loaded elsewhere
% Options for packages loaded elsewhere
\PassOptionsToPackage{unicode}{hyperref}
\PassOptionsToPackage{hyphens}{url}
\PassOptionsToPackage{dvipsnames,svgnames,x11names}{xcolor}
%
\documentclass[
  letterpaper,
  DIV=11,
  numbers=noendperiod]{scrartcl}
\usepackage{xcolor}
\usepackage{amsmath,amssymb}
\setcounter{secnumdepth}{-\maxdimen} % remove section numbering
\usepackage{iftex}
\ifPDFTeX
  \usepackage[T1]{fontenc}
  \usepackage[utf8]{inputenc}
  \usepackage{textcomp} % provide euro and other symbols
\else % if luatex or xetex
  \usepackage{unicode-math} % this also loads fontspec
  \defaultfontfeatures{Scale=MatchLowercase}
  \defaultfontfeatures[\rmfamily]{Ligatures=TeX,Scale=1}
\fi
\usepackage{lmodern}
\ifPDFTeX\else
  % xetex/luatex font selection
\fi
% Use upquote if available, for straight quotes in verbatim environments
\IfFileExists{upquote.sty}{\usepackage{upquote}}{}
\IfFileExists{microtype.sty}{% use microtype if available
  \usepackage[]{microtype}
  \UseMicrotypeSet[protrusion]{basicmath} % disable protrusion for tt fonts
}{}
\makeatletter
\@ifundefined{KOMAClassName}{% if non-KOMA class
  \IfFileExists{parskip.sty}{%
    \usepackage{parskip}
  }{% else
    \setlength{\parindent}{0pt}
    \setlength{\parskip}{6pt plus 2pt minus 1pt}}
}{% if KOMA class
  \KOMAoptions{parskip=half}}
\makeatother
% Make \paragraph and \subparagraph free-standing
\makeatletter
\ifx\paragraph\undefined\else
  \let\oldparagraph\paragraph
  \renewcommand{\paragraph}{
    \@ifstar
      \xxxParagraphStar
      \xxxParagraphNoStar
  }
  \newcommand{\xxxParagraphStar}[1]{\oldparagraph*{#1}\mbox{}}
  \newcommand{\xxxParagraphNoStar}[1]{\oldparagraph{#1}\mbox{}}
\fi
\ifx\subparagraph\undefined\else
  \let\oldsubparagraph\subparagraph
  \renewcommand{\subparagraph}{
    \@ifstar
      \xxxSubParagraphStar
      \xxxSubParagraphNoStar
  }
  \newcommand{\xxxSubParagraphStar}[1]{\oldsubparagraph*{#1}\mbox{}}
  \newcommand{\xxxSubParagraphNoStar}[1]{\oldsubparagraph{#1}\mbox{}}
\fi
\makeatother


\usepackage{longtable,booktabs,array}
\usepackage{calc} % for calculating minipage widths
% Correct order of tables after \paragraph or \subparagraph
\usepackage{etoolbox}
\makeatletter
\patchcmd\longtable{\par}{\if@noskipsec\mbox{}\fi\par}{}{}
\makeatother
% Allow footnotes in longtable head/foot
\IfFileExists{footnotehyper.sty}{\usepackage{footnotehyper}}{\usepackage{footnote}}
\makesavenoteenv{longtable}
\usepackage{graphicx}
\makeatletter
\newsavebox\pandoc@box
\newcommand*\pandocbounded[1]{% scales image to fit in text height/width
  \sbox\pandoc@box{#1}%
  \Gscale@div\@tempa{\textheight}{\dimexpr\ht\pandoc@box+\dp\pandoc@box\relax}%
  \Gscale@div\@tempb{\linewidth}{\wd\pandoc@box}%
  \ifdim\@tempb\p@<\@tempa\p@\let\@tempa\@tempb\fi% select the smaller of both
  \ifdim\@tempa\p@<\p@\scalebox{\@tempa}{\usebox\pandoc@box}%
  \else\usebox{\pandoc@box}%
  \fi%
}
% Set default figure placement to htbp
\def\fps@figure{htbp}
\makeatother





\setlength{\emergencystretch}{3em} % prevent overfull lines

\providecommand{\tightlist}{%
  \setlength{\itemsep}{0pt}\setlength{\parskip}{0pt}}



 
\usepackage[style=alphabetic,]{biblatex}
\addbibresource{references.bib}
\addbibresource{references2.bib}
\addbibresource{google.bib}
\addbibresource{microsoft.bib}
\addbibresource{openai.bib}
\addbibresource{detection.bib}
\addbibresource{ilya.bib}
\addbibresource{creative.bib}
\addbibresource{code.bib}


\KOMAoption{captions}{tableheading}
\makeatletter
\@ifpackageloaded{caption}{}{\usepackage{caption}}
\AtBeginDocument{%
\ifdefined\contentsname
  \renewcommand*\contentsname{Table of contents}
\else
  \newcommand\contentsname{Table of contents}
\fi
\ifdefined\listfigurename
  \renewcommand*\listfigurename{List of Figures}
\else
  \newcommand\listfigurename{List of Figures}
\fi
\ifdefined\listtablename
  \renewcommand*\listtablename{List of Tables}
\else
  \newcommand\listtablename{List of Tables}
\fi
\ifdefined\figurename
  \renewcommand*\figurename{Figure}
\else
  \newcommand\figurename{Figure}
\fi
\ifdefined\tablename
  \renewcommand*\tablename{Table}
\else
  \newcommand\tablename{Table}
\fi
}
\@ifpackageloaded{float}{}{\usepackage{float}}
\floatstyle{ruled}
\@ifundefined{c@chapter}{\newfloat{codelisting}{h}{lop}}{\newfloat{codelisting}{h}{lop}[chapter]}
\floatname{codelisting}{Listing}
\newcommand*\listoflistings{\listof{codelisting}{List of Listings}}
\makeatother
\makeatletter
\makeatother
\makeatletter
\@ifpackageloaded{caption}{}{\usepackage{caption}}
\@ifpackageloaded{subcaption}{}{\usepackage{subcaption}}
\makeatother
\usepackage{bookmark}
\IfFileExists{xurl.sty}{\usepackage{xurl}}{} % add URL line breaks if available
\urlstyle{same}
\hypersetup{
  pdftitle={Fiction Writing Inspired Injection Attacks in LLMs},
  colorlinks=true,
  linkcolor={blue},
  filecolor={Maroon},
  citecolor={blue},
  urlcolor={blue},
  pdfcreator={LaTeX via pandoc}}


\title{Fiction Writing Inspired Injection Attacks in LLMs}
\author{Ana Clara Zoppi Serpa}
\date{}
\begin{document}
\maketitle


\subsection{Literature Review: Open Problems and Literature
Gaps}\label{literature-review-open-problems-and-literature-gaps}

\subsubsection{\texorpdfstring{\autocite{ChenPS025_StruQ}}{{[}@ChenPS025\_StruQ{]}}}\label{chenps025_struq}

\textbf{Paper summary}:

\begin{itemize}
\tightlist
\item
  Attacks rely on LLM ability to follow instructions + inability to
  separate prompts from user data
\item
  Structured queries: general approach to solve the problem. Components:

  \begin{itemize}
  \tightlist
  \item
    \begin{enumerate}
    \def\labelenumi{(\arabic{enumi})}
    \tightlist
    \item
      Secure front-end
    \end{enumerate}
  \item
    \begin{enumerate}
    \def\labelenumi{(\arabic{enumi})}
    \setcounter{enumi}{1}
    \tightlist
    \item
      Specially trained LLM -\textgreater{} \emph{structured instruction
      tuned LLM}
    \end{enumerate}
  \end{itemize}
\item
  Code available in the paper! :)
\item
  Training process:

  \begin{itemize}
  \tightlist
  \item
    Augment standard instruction tuning dataset so that examples have
    instructions in the \textbf{data} part, train the LLM to ignore
    these instructions as they are not in the \textbf{prompt} part.
  \item
    Train new LLM entirely from scratch - infeasible. Structured
    instruction tuning can be applied to \textbf{base LLMs}, so it's not
    from scratch, it's a new form of instruction tuning.
  \end{itemize}
\item
  Special separators, reserved tokens, front-end filters out these
  tokens from user prompt to avoid faking them
\item
  Initialize embedding for the special token before tuning is important
  to keep model performance
\end{itemize}

\textbf{Suggested future work}:

\begin{itemize}
\tightlist
\item
  StruQ decreases ASR of manual attacks tpo 0\% but some optimized
  attacks (TAP, GCG) still have high ASR
\item
  StruQ is trained on task-agnostic manual injections, doesn't fully
  generalize to task-specific injections

  \begin{itemize}
  \tightlist
  \item
    Possible solution: craft task-dependent injections dataset with LLMs
    or human crowdsourcing
  \end{itemize}
\item
  StruQ protects programmatic applications that use API/libraries, but
  web chat-bot users wouldn't like manually marking their inputs as
  prompt x data constantly
\item
  Prompt injection focus: StruQ doesn't prevent jailbreak, data
  extraction or other attacks, just prompt injection
\end{itemize}

\subsubsection{\texorpdfstring{\autocite{GongLZRCC0L25_Papillon}}{{[}@GongLZRCC0L25\_Papillon{]}}}\label{gonglzrcc0l25_papillon}

\textbf{Summary}:

\begin{itemize}
\tightlist
\item
  Automated black-box jailbreaking framework inspired by SWE fuzzy
  testing concepts
\item
  Proprietary APIs e.g.~Gemini, GPT
\item
  Code available
\item
  Differential:

  \begin{itemize}
  \tightlist
  \item
    This framework can start the jailbreaking process from an empty
    seed, no need for starter templates
  \item
    Aims at producing natural, semantically coherent prompts, to bypass
    detection
  \item
    Aims to shorten prompts -\textgreater{} attacks get cheaper because
    proprietary APIs charge on a token basis
  \end{itemize}
\item
  White-box jailbreaking x black-box jailbreaking distinction
\item
  White-box well-known attack: GCG
\item
  Black-box subdivides into handcrafted and automatic
\item
  Jailbreaking defense types

  \begin{itemize}
  \tightlist
  \item
    Input modification
  \item
    Prompt engineering
  \item
    Output filtering
  \end{itemize}
\item
  Defenses Papillon breaks

  \begin{itemize}
  \tightlist
  \item
    perplexity
  \item
    SmoothLLM
  \item
    LlamaGuard
  \end{itemize}
\end{itemize}

\textbf{Suggested future work}:

\begin{itemize}
\tightlist
\item
  Add more mutators to the framework for prompt evolution
\item
  Evaluate different prompt compression methods
\item
  Develop defenses
\item
  Test Papillon on non-English LLM
\end{itemize}

\subsubsection{\texorpdfstring{\autocite{PasquiniKA25_LLMmap}}{{[}@PasquiniKA25\_LLMmap{]}}}\label{pasquinika25_llmmap}

\textbf{Summary}:

\begin{itemize}
\tightlist
\item
  Fingerprinting LLM = figuring out what LLM is being used by an app
  -\textgreater{} analogous to OS fingerprinting
\item
  How it works

  \begin{itemize}
  \tightlist
  \item
    Send carefully crafted queries to the app
  \item
    Analyze responses
  \item
    Identify specific LLM version in use
  \end{itemize}
\item
  42 different LLM versions 95\% accuracy just 8 interactions
\item
  Why is it an attack?

  \begin{itemize}
  \tightlist
  \item
    Fingerprinting enables other attack surfaces, it's related to the
    \textbf{reconnaissance} stage of an attack
  \item
    Knowing the model enables trying target adversarial prompts,
    vulnerabilities specific to that model etc
  \end{itemize}
\item
  Tool for AI red teams
\item
  Code available
\item
  Queries that trigger unique response patterns that might identify what
  LLM is responding

  \begin{itemize}
  \tightlist
  \item
    Meta information (``what is your version?'')
  \item
    Banner grabbing (``what is your name?'')
  \item
    Malformed
  \item
    Alignment / harmful / controversial
  \item
    Nonsensical, semantically broken
  \end{itemize}
\item
  Not necessarily the model answers truthfully to a banner grabbing, but
  it might have a \emph{unique way to deflect the question}
\item
  Collect traces from the queries \& responses
\item
  Use an ML inference model (handle diverse/variable responses)
\item
  Open-set x closed-set

  \begin{itemize}
  \tightlist
  \item
    Closed-set: classifier
  \item
    Open-set: siamese-like, contrastive learning (``is this LLM the same
    as this other LLM?'')
  \end{itemize}
\item
  Prompt configuration universe -\textgreater{} ways LLMs get customized
  and become unique for a certain app

  \begin{itemize}
  \tightlist
  \item
    Hyperparameters (temperature, frequency penalty)
  \item
    System prompt
  \item
    Framework (RAG, COT, other\ldots)
  \end{itemize}
\item
  Some defenses do reduce the accuracy of fingerprinting, but they have
  disadvantages, such as impacting model utility
\item
  Findings suggest fingerprinting is an inevitable consequence of unique
  behavior exhibited by different models, unlike OS fingerprinting,
  which is avoidable with some measures
\end{itemize}

\textbf{Suggested future work}:

\begin{itemize}
\tightlist
\item
  GlitchTokens
\item
  Automated query generation with optimization heuristics
\item
  Additional inference - not only model + version, but also

  \begin{itemize}
  \tightlist
  \item
    tool call list
  \item
    RAG / other internal framework being used
  \item
    hyperparameters
  \end{itemize}
\end{itemize}

\subsubsection{\texorpdfstring{\autocite{BiscontiPRGS26_AdversarialPoetry}}{{[}@BiscontiPRGS26\_AdversarialPoetry{]}}}\label{biscontiprgs26_adversarialpoetry}

\textcolor{red}{to-do - lido no tablet}

\subsubsection{\texorpdfstring{\autocite{BiscontiGPSSN26_AdversarialTales}}{{[}@BiscontiGPSSN26\_AdversarialTales{]}}}\label{biscontigpssn26_adversarialtales}

\textcolor{red}{to-do - lido no tablet}

\subsubsection{\texorpdfstring{\autocite{IBM2025_PromptInjection}}{{[}@IBM2025\_PromptInjection{]}}}\label{ibm2025_promptinjection}

A IBM web page explaining the basics of prompt injection - what it is,
how it generally works, risks etc. Not too technical, not a research
paper with experimental results.

\subsubsection{\texorpdfstring{\autocite{Alber2025_MedicalLLMsAreVulnerableToDataPoisoningAttacks}}{{[}@Alber2025\_MedicalLLMsAreVulnerableToDataPoisoningAttacks{]}}}\label{alber2025_medicalllmsarevulnerabletodatapoisoningattacks}

\textcolor{red}{to-do - tem detalhes nos slides que apresentei no Horus}

\subsubsection{\texorpdfstring{\autocite{Youtube2024_SomeoneWon50kByMakingAIHallucinate},\autocite{IlPost2024_Freysa}}{{[}@Youtube2024\_SomeoneWon50kByMakingAIHallucinate{]},{[}@IlPost2024\_Freysa{]}}}\label{youtube2024_someonewon50kbymakingaihallucinateilpost2024_freysa}

Video and news article about someone winning \$50k by making AI
hallucinate on the Freysa competition.

\subsubsection{\texorpdfstring{\autocite{HeidingKLKS24_LLMPhishing}}{{[}@HeidingKLKS24\_LLMPhishing{]}}}\label{heidingklks24_llmphishing}

\textcolor{red}{não cheguei a ler, tava na lista antiga do google docs}

\subsubsection{\texorpdfstring{\autocite{Kiela2021_TheHatefulMemesChallengeDetectingHateSpeechInMultimodalMemes}}{{[}@Kiela2021\_TheHatefulMemesChallengeDetectingHateSpeechInMultimodalMemes{]}}}\label{kiela2021_thehatefulmemeschallengedetectinghatespeechinmultimodalmemes}

\textcolor{red}{acho que eu li sim, deve estar no Drive ou no tablet, só não lembro os detalhes agora}

\subsubsection{\texorpdfstring{\autocite{Fersini2022_SemEval2022Task5MultimediaAutomaticMisogynyIdentification}}{{[}@Fersini2022\_SemEval2022Task5MultimediaAutomaticMisogynyIdentification{]}}}\label{fersini2022_semeval2022task5multimediaautomaticmisogynyidentification}

\textcolor{red}{n li ainda}

\subsubsection{\texorpdfstring{\autocite{Hwang2023_MemeCapCaptioningAndInterpretingMemes}}{{[}@Hwang2023\_MemeCapCaptioningAndInterpretingMemes{]}}}\label{hwang2023_memecapcaptioningandinterpretingmemes}

\textcolor{red}{n li ainda}

\subsubsection{\texorpdfstring{\autocite{Cardenuto2023_TheAgeOfSyntheticRealitiesChallengesAndOpportunities}}{{[}@Cardenuto2023\_TheAgeOfSyntheticRealitiesChallengesAndOpportunities{]}}}\label{cardenuto2023_theageofsyntheticrealitieschallengesandopportunities}

\textcolor{red}{li sim, é o paper do Anderson e do João, é um geralzão de tudo que tem rolado de realidade sintética}

\subsubsection{\texorpdfstring{\autocite{Farid2016_PhotoForensics}}{{[}@Farid2016\_PhotoForensics{]}}}\label{farid2016_photoforensics}

\textcolor{red}{n li, é o livro de forense digital que fala das técnicas pré-LLM, Anderson tinha recomendado}

\subsubsection{\texorpdfstring{\autocite{Guo2025_DarkMindLatentChainOfThoughtBackdoor}}{{[}@Guo2025\_DarkMindLatentChainOfThoughtBackdoor{]}}}\label{guo2025_darkmindlatentchainofthoughtbackdoor}

\textcolor{red}{to-do - li sim, tem detalhes em slide e no google docs, dá pra pegar}

\subsubsection{Some papers about KG + LLM, didn't
read}\label{some-papers-about-kg-llm-didnt-read}

\autocite{Ibrahim2026},
\autocite{khorashadizadeh2024researchtrendsinterplaylarge},
\autocite{dai2025largelanguagemodelsbetter},
\autocite{kau2024combiningknowledgegraphslarge},
\autocite{Pan2024_UnifyingLargeLanguageModelsAndKnowledgeGraphsARoadmap}

\subsubsection{\texorpdfstring{\autocite{Ng2017_DeepLearningSpecialization}}{{[}@Ng2017\_DeepLearningSpecialization{]}}}\label{ng2017_deeplearningspecialization}

\textcolor{red}{Curso de Deep Learning, relevante, gastei um tempão fazendo pra me especializar}

\subsubsection{Papers about anxiety in
LLMs}\label{papers-about-anxiety-in-llms}

\autocite{codaforno2024inducinganxietylargelanguage}
\autocite{Ben-Zion2025}

\subsubsection{Milestone papers about
LLMs}\label{milestone-papers-about-llms}

These are really cool to see the historical trajectory, taken from this
repository with several LLM resources and papers:
\autocite{Hannibal046AwesomeLLM}

\begin{itemize}
\tightlist
\item
  \autocite{vaswani2023attentionneed} - Attention Is All You Need
\item
  \autocite{Radford2018ImprovingLU} - Improving Language Understanding
  by Generative Pre-Training
\item
  \autocite{devlin2019bertpretrainingdeepbidirectional} - BERT:
  Pre-training of Deep Bidirectional Transformers for Language
  Understanding
\item
  \autocite{Radford2019LanguageMA} - Language Models are Unsupervised
  Multitask Learners
\item
  \autocite{shoeybi2020megatronlmtrainingmultibillionparameter} -
  Megatron-LM: Training Multi-Billion Parameter Language Models Using
  Model Parallelism
\item
  \autocite{JMLR:v21:20-074} - Exploring the Limits of Transfer Learning
  with a Unified Text-to-Text Transformer
\item
  \autocite{rajbhandari2020zeromemoryoptimizationstraining} - ZeRO:
  Memory Optimizations Toward Training Trillion Parameter Models
\item
  \autocite{kaplan2020scalinglawsneurallanguage} - Scaling Laws for
  Neural Language Models
\item
  \autocite{brown2020languagemodelsfewshotlearners} - Language Models
  are Few-Shot Learners
\item
  \autocite{fedus2022switchtransformersscalingtrillion} - Switch
  Transformers: Scaling to Trillion Parameter Models with Simple and
  Efficient Sparsity
\item
  \autocite{chen2021evaluatinglargelanguagemodels} - Evaluating Large
  Language Models Trained on Code
\item
  \autocite{bommasani2022opportunitiesrisksfoundationmodels} - On the
  Opportunities and Risks of Foundation Models
\item
  \autocite{wei2022finetunedlanguagemodelszeroshot} - Finetuned Language
  Models Are Zero-Shot Learners
\item
  \autocite{sanh2022multitaskpromptedtrainingenables} - Multitask
  Prompted Training Enables Zero-Shot Task Generalization
\item
  \autocite{du2022glamefficientscalinglanguage} - GLaM: Efficient
  Scaling of Language Models with Mixture-of-Experts
\item
  \autocite{Nakano2021WebGPTBQ} - WebGPT: Browser-assisted
  question-answering with human feedback
\item
  \autocite{borgeaud2022improvinglanguagemodelsretrieving} - Improving
  language models by retrieving from trillions of tokens
\item
  \autocite{rae2022scalinglanguagemodelsmethods} - Scaling Language
  Models: Methods, Analysis \& Insights from Training Gopher
\item
  \autocite{wei2023chainofthoughtpromptingelicitsreasoning} -
  Chain-of-Thought Prompting Elicits Reasoning in Large Language Models
\item
  \autocite{thoppilan2022lamdalanguagemodelsdialog} - LaMDA: Language
  Models for Dialog Applications
\item
  \autocite{lewkowycz2022solvingquantitativereasoningproblems} - Solving
  Quantitative Reasoning Problems with Language Models
\item
  \autocite{smith2022usingdeepspeedmegatrontrain} - Using DeepSpeed and
  Megatron to Train Megatron-Turing NLG 530B, A Large-Scale Generative
  Language Model
\item
  \autocite{ouyang2022traininglanguagemodelsfollow} - Training language
  models to follow instructions with human feedback
\item
  \autocite{chowdhery2022palmscalinglanguagemodeling} - PaLM: Scaling
  Language Modeling with Pathways
\item
  \autocite{hoffmann2022trainingcomputeoptimallargelanguage} - Training
  Compute-Optimal Large Language Models
\item
  \autocite{zhang2022optopenpretrainedtransformer} - OPT: Open
  Pre-trained Transformer Language Models
\item
  \autocite{tay2023ul2unifyinglanguagelearning} - UL2: Unifying Language
  Learning Paradigms
\item
  \autocite{wei2022emergentabilitieslargelanguage} - Emergent Abilities
  of Large Language Models
\item
  \autocite{srivastava2023beyond} - Beyond the Imitation Game:
  Quantifying and extrapolating the capabilities of language models
\item
  \autocite{hao2022languagemodelsgeneralpurposeinterfaces} - Language
  Models are General-Purpose Interfaces
\item
  \autocite{glaese2022improvingalignmentdialogueagents} - Improving
  alignment of dialogue agents via targeted human judgements
\item
  \autocite{chung2022scalinginstructionfinetunedlanguagemodels} -
  Scaling Instruction-Finetuned Language Models
\item
  \autocite{zeng2023glm130bopenbilingualpretrained} - GLM-130B: An Open
  Bilingual Pre-trained Model
\item
  \autocite{liang2023holisticevaluationlanguagemodels} - Holistic
  Evaluation of Language Models
\item
  \autocite{workshop2023bloom176bparameteropenaccessmultilingual} -
  BLOOM: A 176B-Parameter Open-Access Multilingual Language Model
\item
  \autocite{taylor2022galacticalargelanguagemodel} - Galactica: A Large
  Language Model for Science
\item
  \autocite{iyer2023optimlscalinglanguagemodel} - OPT-IML: Scaling
  Language Model Instruction Meta Learning through the Lens of
  Generalization
\item
  \autocite{longpre2023flancollectiondesigningdata} - The Flan
  Collection: Designing Data and Methods for Effective Instruction
  Tuning
\item
  \autocite{touvron2023llamaopenefficientfoundation} - LLaMA: Open and
  Efficient Foundation Language Models
\item
  \autocite{huang2023languageneedaligningperception} - Language Is Not
  All You Need: Aligning Perception with Language Models
\item
  \autocite{orvieto2023resurrectingrecurrentneuralnetworks} -
  Resurrecting Recurrent Neural Networks for Long Sequences
\item
  \autocite{driess2023palmeembodiedmultimodallanguage} - PaLM-E: An
  Embodied Multimodal Language Model
\item
  \autocite{openai2024gpt4technicalreport} - GPT-4 Technical Report
\item
  \autocite{liu2023visualinstructiontuning} - Visual Instruction Tuning
\item
  \autocite{biderman2023pythiasuiteanalyzinglarge} - Pythia: A Suite for
  Analyzing Large Language Models Across Training and Scaling
\item
  \autocite{sun2023principledrivenselfalignmentlanguagemodels} -
  Principle-Driven Self-Alignment of Language Models from Scratch with
  Minimal Human Supervision
\item
  \autocite{anil2023palm2technicalreport} - PaLM 2 Technical Report
\item
  \autocite{peng2023rwkvreinventingrnnstransformer} - RWKV: Reinventing
  RNNs for the Transformer Era
\item
  \autocite{rafailov2024directpreferenceoptimizationlanguage} - Direct
  Preference Optimization: Your Language Model is Secretly a Reward
  Model
\item
  \autocite{yao2023treethoughtsdeliberateproblem} - Tree of Thoughts:
  Deliberate Problem Solving with Large Language Models
\item
  \autocite{touvron2023llama2openfoundation} - Llama 2: Open Foundation
  and Fine-Tuned Chat Models
\item
  \autocite{jiang2023mistral7b} - Mistral 7B
\item
  \autocite{gu2024mambalineartimesequencemodeling} - Mamba: Linear-Time
  Sequence Modeling with Selective State Spaces
\item
  \autocite{deepseekai2024deepseekv2strongeconomicalefficient} -
  DeepSeek-V2: A Strong, Economical, and Efficient Mixture-of-Experts
  Language Model
\item
  \autocite{groeneveld-etal-2024-olmo} - OLMo: Accelerating the Science
  of Language Models
\item
  \autocite{dao2024transformersssmsgeneralizedmodels} - Transformers are
  SSMs: Generalized Models and Efficient Algorithms Through Structured
  State Space Duality
\item
  \autocite{grattafiori2024llama3herdmodels} - The Llama 3 Herd of
  Models
\item
  \autocite{penedo2024finewebdatasetsdecantingweb} - The FineWeb
  Datasets: Decanting the Web for the Finest Text Data at Scale
\item
  \autocite{muennighoff2025olmoeopenmixtureofexpertslanguage} - OLMoE:
  Open Mixture-of-Experts Language Models
\item
  \autocite{qwen2025qwen25technicalreport} - Qwen2.5 Technical Report
\item
  \autocite{deepseekai2025deepseekv3technicalreport} - DeepSeek-V3
  Technical Report
\item
  \autocite{Guo_2025} - DeepSeek-R1 incentivizes reasoning in LLMs
  through reinforcement learning
\end{itemize}

\subsubsection{Google Research}\label{google-research}

\begin{itemize}
\tightlist
\item
  \autocite{52868} - InstructPipe: Generating Visual Blocks Pipelines
  with Human Instructions and LLMs
\item
  \autocite{1021008} - Gemini \& Physical World: Large Language Models
  Can Estimate the Intensity of Earthquake Shaking from Multi-Modal
  Social Media Posts
\item
  \autocite{1025553} - SSDTrain: Faster Large Language Model Training
  Using SSD-Based Activation Offloading
\item
  \autocite{1024745} - Triaging mammography with artificial
  intelligence: an implementation study
\item
  \autocite{1024143} - Mufu: Multilingual Fused Learning for
  Low-Resource Translation with LLM
\item
  \autocite{1026765} - TOKENFORMER: Rethinking Transformers Scaling with
  Tokenized Model Parameters
\item
  \autocite{morris2024generativeghostsanticipatingbenefits} - Generative
  Ghosts: Anticipating Benefits and Risks of AI Afterlives
\item
  \autocite{lee2025geminiembeddinggeneralizableembeddings} - Gemini
  Embedding: Generalizable Embeddings from Gemini
\item
  \autocite{maninis2025tipstextimagepretrainingspatial} - TIPS:
  Text-Image Pretraining with Spatial awareness
\item
  \autocite{wang2025scalingpretrainingbilliondata} - Scaling
  Pre-training to One Hundred Billion Data for Vision Language Models
\item
  \autocite{cooper2025machineunlearningdoesntthink} - Machine Unlearning
  Doesn't Do What You Think: Lessons for Generative AI Policy and
  Research
\item
  \autocite{schultz2025masteringboardgamesexternal} - Mastering Board
  Games by External and Internal Planning with Language Models
\item
  \autocite{dambrosio2025achievinghumanlevelcompetitive} - Achieving
  Human Level Competitive Robot Table Tennis
\item
  \autocite{Hutter:24uaibook2} - An Introduction to Universal Artificial
  Intelligence
\item
  \autocite{panda2024teach} - Teach LLMs to Phish: Stealing Private
  Information from Language Models
\item
  \autocite{whitney2024learning} - Learning 3D Particle-based Simulators
  from RGB-D Videos
\item
  \autocite{song2024positionleveragefoundationalmodels} - Position:
  Leverage Foundational Models for Black-Box Optimization
\end{itemize}

\subsubsection{Microsoft Research}\label{microsoft-research}

\begin{itemize}
\tightlist
\item
  \autocite{vorvoreanu2025fostering} - Fostering appropriate reliance on
  GenAI: Lessons learned from early research
\item
  \autocite{flores-saviaga2025the} - The Impact of Generative AI Coding
  Assistants on Developers Who Are Visually Impaired
\item
  \autocite{alshammari2025a} - A Unifying Framework for Representation
  Learning
\end{itemize}

\subsubsection{OpenAI}\label{openai}

\begin{itemize}
\tightlist
\item
  \autocite{miserendino2025swelancerfrontierllmsearn} - SWE-Lancer: Can
  Frontier LLMs Earn \$1 Million from Real-World Freelance Software
  Engineering?
\item
  \autocite{baker2022videopretrainingvptlearning} - Video PreTraining
  (VPT): Learning to Act by Watching Unlabeled Online Videos
\item
  \autocite{radford2021learningtransferablevisualmodels} - Learning
  Transferable Visual Models From Natural Language Supervision (CLIP)
\item
  \autocite{openai2021dalle} - DALL·E: Creating Images from Text
\item
  \autocite{openai2018montezumasrevenge} - Learning Montezuma's Revenge
  from a single demonstration
\item
  \autocite{openai2022languagemodelsafety} - Lessons learned on language
  model safety and misuse
\item
  \autocite{goh2021multimodal} - Multimodal Neurons in Artificial Neural
  Networks
\end{itemize}

\subsubsection{Synthetic Media
Detection}\label{synthetic-media-detection}

\begin{itemize}
\tightlist
\item
  \autocite{9141516} - FakeCatcher: Detection of Synthetic Portrait
  Videos using Biological Signals
\item
  \autocite{unknown} - Synthetic Realities in the Digital Age:
  Navigating the Opportunities and Challenges of AI-Generated Content
\item
  \autocite{ye2025lokicomprehensivesyntheticdata} - LOKI: A
  Comprehensive Synthetic Data Detection Benchmark using Large
  Multimodal Models
\item
  \autocite{Zobaed2021} - DeepFakes: Detecting Forged and Synthetic
  Media Content Using Machine Learning
\item
  \autocite{10095167} - On The Detection of Synthetic Images Generated
  by Diffusion Models
\item
  \autocite{10409290} - CIFAKE: Image Classification and Explainable
  Identification of AI-Generated Synthetic Images
\item
  \autocite{laurier2024catmousegameongoing} - The Cat and Mouse Game:
  The Ongoing Arms Race Between Diffusion Models and Detection Methods
\end{itemize}

\subsubsection{Ilya Sutskever's Recommended Reading
List}\label{ilya-sutskevers-recommended-reading-list}

Papers and resources from \autocite{dzyim2023ilyasutskeverreading},
originally shared by Ilya Sutskever with John Carmack in 2020.

\begin{itemize}
\tightlist
\item
  \autocite{rush2018annotated} - The Annotated Transformer
\item
  \autocite{aaronson2011first} - The First Law of Complexodynamics
\item
  \autocite{karpathy2015unreasonable} - The Unreasonable Effectiveness
  of Recurrent Neural Networks
\item
  \autocite{olah2015understanding} - Understanding LSTM Networks
\item
  \autocite{zaremba2014recurrent} - Recurrent Neural Network
  Regularization
\item
  \autocite{hinton1993keeping} - Keeping the Neural Networks Simple by
  Minimizing the Description Length of the Weights
\item
  \autocite{vinyals2015pointer} - Pointer Networks
\item
  \autocite{krizhevsky2012imagenet} - ImageNet Classification with Deep
  Convolutional Neural Networks
\item
  \autocite{vinyals2016order} - Order Matters: Sequence to Sequence for
  Sets
\item
  \autocite{huang2019gpipe} - GPipe: Easy Scaling with Micro-Batch
  Pipeline Parallelism
\item
  \autocite{he2016deep} - Deep Residual Learning for Image Recognition
\item
  \autocite{yu2016multi} - Multi-Scale Context Aggregation by Dilated
  Convolutions
\item
  \autocite{gilmer2017neural} - Neural Message Passing for Quantum
  Chemistry
\item
  \autocite{bahdanau2016neural} - Neural Machine Translation by Jointly
  Learning to Align and Translate
\item
  \autocite{he2016identity} - Identity Mappings in Deep Residual
  Networks
\item
  \autocite{santoro2017simple} - A Simple Neural Network Module for
  Relational Reasoning
\item
  \autocite{chen2017variational} - Variational Lossy Autoencoder
\item
  \autocite{santoro2018relational} - Relational Recurrent Neural
  Networks
\item
  \autocite{aaronson2014quantifying} - Quantifying the Rise and Fall of
  Complexity in Closed Systems: The Coffee Automaton
\item
  \autocite{graves2014neural} - Neural Turing Machines
\item
  \autocite{amodei2016deep} - Deep Speech 2: End-to-End Speech
  Recognition in English and Mandarin
\item
  \autocite{grunwald2004tutorial} - A Tutorial Introduction to the
  Minimum Description Length Principle
\item
  \autocite{legg2008machine} - Machine Super Intelligence (PhD Thesis)
\item
  \autocite{shen2017kolmogorov} - Kolmogorov Complexity and Algorithmic
  Randomness
\item
  \autocite{feifei2016cs231n} - CS231n: Convolutional Neural Networks
  for Visual Recognition
\end{itemize}

\subsubsection{Misc, to be organized}\label{misc-to-be-organized}

\begin{itemize}
\tightlist
\item
  \autocite{kumar-etal-2023-language} - Language Generation Models Can
  Cause Harm: So What Can We Do About It? An Actionable Survey
\item
  \autocite{agrawal2024knowledgegraphsreducehallucinations} - Can
  Knowledge Graphs Reduce Hallucinations in LLMs? : A Survey
\item
  \autocite{feng2023knowledgesolverteachingllms} - Knowledge Solver:
  Teaching LLMs to Search for Domain Knowledge from Knowledge Graphs
\item
  \autocite{lu2024aiscientistfullyautomated} - The AI Scientist: Towards
  Fully Automated Open-Ended Scientific Discovery
\item
  \autocite{doi:10.1126/science.adq1814} - Durably reducing conspiracy
  beliefs through dialogues with AI
\item
  \autocite{nakanishi2025scalablesoftmaxsuperiorattention} -
  Scalable-Softmax Is Superior for Attention
\item
  \autocite{feng2024rnnsneeded} - Were RNNs All We Needed?
\item
  \autocite{zhang2025thoughtactionhierarchyneural} - From Thought to
  Action: How a Hierarchy of Neural Dynamics Supports Language
  Production
\item
  \autocite{lévy2025braintotextdecodingnoninvasiveapproach} -
  Brain-to-Text Decoding: A Non-invasive Approach via Typing
\item
  \autocite{li-etal-2025-loki} - Loki: An Open-Source Tool for Fact
  Verification
\end{itemize}


\printbibliography



\end{document}
